% SGA 3，第 VII 篇_B，英文重建本——第 1 轮（中文译本）。
% 范围：完整包含第 1.3.5 节，始于依据原文件第 29 页；止于原文件第 31 页的 1.3.6 之前。
% 合并阅读版页码： 542--544. 再版印刷页码： 531--533.
% 所译正文占原文件第 29--30 页；原文件第 31 页确定
% 终止边界。 原版页边定位： 502--503.
% 本范围内的编者注： 60--61.
% 依据： Polo--Gille Exp7B-13oct24.pdf,
% SHA-256 D303690E0195654BBE7C3DAA5D89EA23E24CD709A3575217234AC19AE8836694.
% 曾参阅 Jacob Reinhold 的 CC BY 4.0 Markdown，作为对比版本来源。
% 本范围内无原生图；所有陈列数学均为
% 使用原生 LaTeX 重建。

\subsubsection*{1.3.5}
\phantomsection
\label{sga3:VIIB:1.3.5}

设 \(k\) 是伪紧环。若 \(k\) 上的形式簇 \(X\) 的仿射代数 \(A\) 是
\emph{投射}伪紧 \(k\)-模，即结构态射
\[
  X\longrightarrow\operatorname{Spf}(k)
\]
拓扑平坦，则称这个形式簇是\emph{拓扑平坦}的。

首先注意，\footnote{\label{sga3:VIIB:nde:60}编者添加了原文中隐含使用的下列引理，
并引入编号 \SGARefLink{sga3:VIIB:1.3.5.A}{1.3.5.A}--\SGARefLink{sga3:VIIB:1.3.5.D}{1.3.5.D}，以供下文引用。}
\SGARefLink{sga3:VIIB:0.2.2}{0.2.2} 和 \SGARefLink{sga3:VIIB:0.3.6}{0.3.6}
给出下列结果，它与 VII\textsubscript{A}，\SGARefLink{sga3:VIIA:3.1.1}{3.1.1} 类似。

\medskip
\noindent\textbf{引理 1.3.5.A。}
\phantomsection
\label{sga3:VIIB:1.3.5.A}
{\itshape 假定 \(k\) 是 Artin 环。函子
\[
  A\longmapsto A^\dagger=\operatorname{Hom}_c(A,k),
  \qquad
  C\longmapsto C^*=\operatorname{Hom}_k(C,k)
\]
在拓扑平坦的预有限 \(k\)-代数范畴与平坦 \(k\)-余代数范畴之间定义反等价。}

事实上，若 \(A\) 是拓扑平坦的预有限 \(k\)-代数，则由
\SGARefLink{sga3:VIIB:0.3.6}{0.3.6}，有 \(k\)-模同构
\[
  A^\dagger\otimes_k A^\dagger
  \xrightarrow{\ \sim\ }
  \bigl(A\widehat{\otimes}A\bigr)^\dagger.
\]
所以乘法
\[
  A\widehat{\otimes}A\longrightarrow A
\]
通过对偶性在 \(A^\dagger\) 上诱导一个 \(k\)-余代数结构。
其余结论由\SGARefLink{sga3:VIIB:0.2.2}{命题 0.2.2} 得到。

回到任意伪紧环 \(k\) 的情形。

\medskip
\noindent\textbf{定义 1.3.5.B。}
\phantomsection
\label{sga3:VIIB:1.3.5.B}
对每个仿射代数 \(A\) 是投射伪紧 \(k\)-模的形式 \(k\)-簇 \(X\)，
可以对应一个平坦 \(\mathbf O_k\)-余代数 \(\mathbf H(X)\)，定义如下。

用 \(\mathbf H(X)\) 表示 \(\mathbf O_k\)-模 \(\mathbf V_k^f(A)\)，即“\(A\) 的对偶”。
由于 \(A\) 的底伪紧 \(k\)-模是投射模，它是平坦 \(\mathbf O_k\)-模
\textup{（参见\ \SGARefLink{sga3:VIIB:1.2.3.C}{1.2.3.C}）}。此外，由
% 依据：原文件第 30 页；合并阅读版第 543 页；印刷页码 532。
\SGARefLink{sga3:VIIB:0.3.6}{0.3.6}，有同构
\[
  \mathbf V_k^f\bigl(A\widehat{\otimes}A\bigr)
  \simeq
  \mathbf V_k^f(A)\otimes\mathbf V_k^f(A).
\]
因此，\(A\) 的乘法通过转置诱导对角态射
\[
  \mathbf H(X)=\mathbf V_k^f(A)
  \longrightarrow
  \mathbf V_k^f\bigl(A\widehat{\otimes}A\bigr)
  =\mathbf H(X)\otimes\mathbf H(X),
\]
它使 \(\mathbf H(X)\) 成为平坦 \(\mathbf O_k\)-余代数。
我们称 \(\mathbf H(X)\) 为 \emph{\(X\) 在 \(\mathbf O_k\) 上的余代数}。

\medskip
\noindent\textbf{定义 1.3.5.C。}
\phantomsection
\label{sga3:VIIB:1.3.5.C}
反之，对每个 \(\mathbf O_k\)-余代数 \(\mathbf C\)，可以对应一个 \(k\)-函子
\textup{（参见\ \SGARefLink{sga3:VIIA:1.2.1}{1.2.1}）}\(\operatorname{Spf}^*(\mathbf C)\)，
定义如下。对 \(\operatorname{Alf}/k\) 的每个对象 \(B\)，使用
VII\textsubscript{A} \SGARefLink{sga3:VIIA:3.1}{3.1} 的记号，令
\begin{equation}
\begin{aligned}
  \operatorname{Spf}^*(\mathbf C)(B)
  &=
  \operatorname{Hom}_{B\text{-余代数}}\bigl(B,\mathbf C(B)\bigr)
  \\
  &=
  \bigl\{x\in\mathbf C(B)\mid
    \varepsilon_{\mathbf C(B)}(x)=1
    \ \text{且}\
    \Delta_{\mathbf C(B)}(x)=x\otimes x\bigr\}.
\end{aligned}
\tag{1}
\label{sga3:VIIB:1.3.5.C:1}
\end{equation}

再假定\footnote{\label{sga3:VIIB:nde:61}编者补充了以下细节。}
\(\mathbf C\) 的底 \(\mathbf O_k\)-模是容许的
\textup{（参见\ \SGARefLink{sga3:VIIA:1.2.1}{1.2.1}）}，并令
\begin{equation}
  A
  =
  \Gamma^*(\mathbf C)
  =
  \varprojlim_{\mathfrak l}\mathbf C(k/\mathfrak l)^*.
\tag{2}
\label{sga3:VIIB:1.3.5.C:2}
\end{equation}
各模 \(\mathbf C(k/\mathfrak l)^*\) 上的代数结构使 \(A\) 配备一个预有限
\(k\)-代数结构。对 \(\operatorname{Alf}/k\) 的每个对象 \(B\)，有
\begin{equation}
\begin{aligned}
  \operatorname{Hom}_{\operatorname{Vaf}/k}
    \bigl(\operatorname{Spf}(B),\operatorname{Spf}(A)\bigr)
  &=
  \operatorname{Hom}_{\operatorname{Alp}/k}(A,B)
  \\
  &=
  \operatorname{Hom}_{\operatorname{Alp}/B}
    \bigl(A\widehat{\otimes}B,B\bigr).
\end{aligned}
\tag{3}
\label{sga3:VIIB:1.3.5.C:3}
\end{equation}

最后假定 \(\mathbf C\) 是平坦 \(\mathbf O_k\)-模。
则由 \SGARefLink{sga3:VIIB:1.2.3.E}{1.2.3.E}，\(A=\Gamma^*(\mathbf C)\)
是投射伪紧 \(k\)-模。此外，我们在\emph{同处}的证明中已经看到，
若 \(\mathfrak l_0\) 是 \(k\) 中包含于 \(k\to B\) 的核的开理想，则有同构
\begin{equation}
\begin{aligned}
  A\widehat{\otimes}B
  &=
  \varprojlim_{\mathfrak l}
    \mathbf C(k/\mathfrak l)^*\widehat{\otimes}_kB
  \\
  &\simeq
  \operatorname{Hom}_{k/\mathfrak l_0}
    \bigl(\mathbf C(k/\mathfrak l_0),B\bigr)
  \\
  &\simeq
  \operatorname{Hom}_B\bigl(\mathbf C(B),B\bigr).
\end{aligned}
\tag{4}
\label{sga3:VIIB:1.3.5.C:4}
\end{equation}
我们将右端记为 \(\mathbf C(B)^*\)。最后，将
\SGARefLink{sga3:VIIB:1.3.5.A}{引理 1.3.5.A} 应用于 Artin 环 \(B\)，得到自然同构
\begin{equation}
  \operatorname{Hom}_{B\text{-余代数}}\bigl(B,\mathbf C(B)\bigr)
  \xrightarrow{\ \sim\ }
  \operatorname{Hom}_{\operatorname{Alp}/B}
    \bigl(\mathbf C(B)^*,B\bigr).
\tag{5}
\label{sga3:VIIB:1.3.5.C:5}
\end{equation}
结合 \textup{\SGARefLink{sga3:VIIB:1.3.5.C:1}{(1)}}、
\textup{\SGARefLink{sga3:VIIB:1.3.5.C:5}{(5)}}、
\textup{\SGARefLink{sga3:VIIB:1.3.5.C:4}{(4)}}、
\textup{\SGARefLink{sga3:VIIB:1.3.5.C:3}{(3)}} 和 \textup{(2)}，
在 \(\mathbf C\) 是平坦 \(\mathbf O_k\)-余代数时，得到函子同构
\begin{equation}
  \operatorname{Spf}^*(\mathbf C)
  \simeq
  \operatorname{Spf}(A)
  =
  \operatorname{Spf}\bigl(\Gamma^*(\mathbf C)\bigr).
\tag{\(\star\)}
\label{sga3:VIIB:1.3.5.C:star}
\end{equation}

因此，若 \(\mathcal A(X)\) 表示形式 \(k\)-簇 \(X\) 的仿射代数，
则结合 \SGARefLink{sga3:VIIB:1.2.3.E}{1.2.3.E}，得到如下结论。

\medskip
\noindent\textbf{命题 1.3.5.D。}
\phantomsection
\label{sga3:VIIB:1.3.5.D}
{\itshape
\begin{enumerate}
  \item[\textup{(i)}]
  函子
  \[
    X\longmapsto
    \mathbf H(X)=\mathbf V_k^f\bigl(\mathcal A(X)\bigr),
    \qquad
    \mathbf C\longmapsto
    \operatorname{Spf}^*(\mathbf C)
    =
    \operatorname{Spf}\bigl(\Gamma^*(\mathbf C)\bigr)
  \]
  在拓扑平坦的形式 \(k\)-簇范畴与平坦 \(\mathbf O_k\)-余代数范畴之间定义等价。

  \item[\textup{(ii)}]
  此外，这个等价“与基变换交换”：若 \(k\to K\) 是伪紧环态射，
  则 \(X\widehat{\otimes}_kK\) 对应于 \(\mathbf H(X)\otimes_kK\)。
\end{enumerate}}

% 下一断点： Exp7B 原文件页码 31, 合并阅读版页码 544,
% 再版印刷页码 533, 紧接其前： 1.3.6.
